\scanpage{7}

\unsure{isons|isom|isom.}, on sait déjà qu'on a \deleted{\unsure{une}} un ouvert \uline{quasi-projectif} sur $S$ du schéma de Hilbert [car le \deleted{\unsure{pro}} degré de $C$ dans $A_T$ est connu \unsure{--|- -}]. \deleted{Il faut} La deuxième condition elle-même s'exprime par un ouvert plus petit.

On a ainsi prouvé que $F'$, \deleted{donc $F$}, est représentable quasi-projectivement, donc $F$ aussi. Nous voulons prouver que l'on a un \deleted{\uline{immersion}} \inserted{morphisme radiciel etc}. Par le critère \unsure{valuatif|habituel}, cela nous ramène aux \deleted{\unsure{de|deux cas}} particuliers suivants :

a) $S=\mathrm{Spec}(k)$ spectre d'un corps, alg. clos si on veut. a) \unsure{$\exists$}$M$ ne peut avoir que \deleted{\unsure{un}} \inserted{\unsure{au plus} deux} pt \unsure{ent $k$|sur $k$}, \inserted{conjugués par le groupe $\mathbb Z/2\subset\mathrm{Aut}_{\unsure{\mathrm{gr}|S}}(A)$} : c'est le \uline{théorème de Torelli} sous sa forme classique

b) \inserted{En \unsure{un|tout|fait} pt $x\in M$ autre que} $M\longrightarrow\mathrm{Spec}(k)$ est \uline{non ramifié}, i.e. \deleted{tt} \inserted{deux} pts de $M$ à valeurs ds $\mathrm{Spec}(k[\epsilon]/\epsilon^2)$ \deleted{\unsure{et ont pour}} \inserted{localisés en le même point de $M$} sont identiques. \deleted{\unsure{P|Pour}} Ce dernier point,

% archive sheet number: 6
